\providecommand{\tightlist}{\setlength{\itemsep}{0pt}\setlength{\parskip}{0pt}}
\newpage
\section{Algoritmos variacionais}
\subsection{Simulação hamiltoniana e trotterização}
Um computador quântico é, antes de tudo, bom em imitar outros sistemas quânticos: uma molécula, um ímã, um material. Imitar é deixar o sistema evoluir no tempo, e isso é calcular \(e^{-iHt}\) para um operador \(H\) (o hamiltoniano). Como transformar essa exponencial em portas de 1 e 2 qubits?

\subsubsection{Evolução no tempo}\label{evoluuxe7uxe3o-no-tempo}

O \textbf{hamiltoniano} \(H\) é o operador hermitiano que representa a energia do sistema. Se o estado inicial é \(\ket{\psi(0)}\), depois de um tempo \(t\) ele vira

\[
\ket{\psi(t)} = e^{-iHt}\ket{\psi(0)}.
\]

A exponencial de uma matriz \(A\) é definida pela mesma série da exponencial comum: \(e^{A} = I + A + A^2/2! + \dots\). Quando \(H\) é hermitiano, \(e^{-iHt}\) é unitária, isto é, é uma porta quântica válida.

Um caso fácil: \(H = Z\). Como \(Z^2 = I\), a série se reorganiza em cosseno e seno:

\[
e^{-i\theta Z} = \cos\theta\, I - i\sin\theta\, Z = \begin{pmatrix} e^{-i\theta} & 0 \\ 0 & e^{i\theta}\end{pmatrix}.
\]

Isso é a rotação \(R_z(2\theta)\) que você já viu, pois \(R_z(\varphi) = e^{-i\varphi Z/2}\). O mesmo vale para \(X\) e \(Y\) (qualquer Pauli \(P\) tem \(P^2=I\)): \(e^{-i\theta P} = \cos\theta\,I - i\sin\theta\,P\).

Qual é o ângulo \(\varphi\) do \(R_z(\varphi)\) equivalente a \(e^{-i\theta Z}\) com \(\theta = 0{,}3\)?

\subsubsection{H como soma de termos de Pauli}\label{h-como-soma-de-termos-de-pauli}

Sistemas reais têm hamiltonianos com vários pedaços. Todo hamiltoniano de \(n\) qubits pode ser escrito como soma de produtos de Paulis (strings de Pauli), cada uma com um coeficiente real:

\[
H = \sum_k c_k\, P_k, \qquad \text{por exemplo } H = 0{,}5\,Z_1 Z_2 + 0{,}3\,X_1 + 0{,}3\,X_2.
\]

Aqui \(Z_1Z_2\) quer dizer \(Z\otimes Z\), \(X_1\) quer dizer \(X\otimes I\), e assim por diante. Cada termo, sozinho, é fácil de exponenciar. O problema: \(e^{-i(A+B)t}\) \textbf{não} é, em geral, \(e^{-iAt}e^{-iBt}\). Isso só vale se \(A\) e \(B\) \textbf{comutam} (\(AB = BA\)).

\(Z_1\) e \(Z_2\) comutam? E \(Z_1\) e \(X_1\)? Responda ``sim'' ou ``não'' para o primeiro par: vale \(e^{-i(Z_1+Z_2)t} = e^{-iZ_1t}e^{-iZ_2t}\)?

\(Z_1\) e \(X_1\) (ambos no qubit 1) não comutam, e é aí que mora o erro.

\subsubsection{Por que o produto simples falha}\label{por-que-o-produto-simples-falha}

Teste com números: \(H = X + Z\), \(t = 1\), estado inicial \(\ket{0}\). A evolução exata dá probabilidade \(P(0) \approx 0{,}512\) de medir 0. Se trocamos \(e^{-i(X+Z)}\) por \(e^{-iX}e^{-iZ}\) (um único ``passo''), obtemos \(P(0) \approx 0{,}292\). A diferença é grande, e vem de \(X\) e \(Z\) não comutarem.

\subsubsection{A fórmula de Trotter}\label{a-fuxf3rmula-de-trotter}

A ideia: dividir o tempo \(t\) em \(n\) passos curtos e, em cada um, aplicar os termos um de cada vez, como se comutassem. Para \(H = A + B\):

\[
e^{-i(A+B)t} \;\approx\; \left( e^{-iAt/n}\, e^{-iBt/n} \right)^{n}.
\]

Essa é a fórmula de Trotter (e Suzuki generalizou o método a ordens maiores). Chama-se de \textbf{1ª ordem} porque o erro de cada passo é de segunda ordem no tamanho do passo, \((t/n)^2\), e é o esquema mais simples. Quanto menor o passo, menos importa a ordem dos termos, e no limite \(n\to\infty\) a fórmula é exata.

Cada letra pode ser um grupo de termos que comutam entre si (assim o grupo se exponencia exatamente). Em \(H = 0{,}5\,Z_1Z_2 + 0{,}3\,X_1 + 0{,}3\,X_2\), agrupe \(A = 0{,}5\,Z_1Z_2\) e \(B = 0{,}3(X_1 + X_2)\): \(X_1\) e \(X_2\) comutam (estão em qubits diferentes), então \(e^{-iBt/n}\) é simplesmente uma rotação \(R_x\) em cada qubit.

Com \(n = 100\) passos, o caso \(H = X+Z\), \(t=1\) acima dá \(P(0)\approx 0{,}5121\), quase o valor exato \(0{,}5122\).

\paragraph{Tamanho do erro}\label{tamanho-do-erro}

Cada passo erra por algo proporcional a \((t/n)^2\) vezes o comutador \([A,B] = AB - BA\). Somando \(n\) passos, um limite típico é

\[
\text{erro} \lesssim \frac{\lVert [A,B]\rVert\, t^2}{2n},
\]

onde \(\lVert\cdot\rVert\) é o tamanho (norma) da matriz. Não é uma lei exata, é um limite que vale quando esse número é pequeno. O que importa é a escala, \(\text{erro} \sim t^2/n\).

Exemplo: com \(t = 1\) e \(n = 10\) o erro é da ordem de \(0{,}1\); dobrando para \(n=20\), cai para \(0{,}05\). Dobrar o tempo simulado quadruplica o erro, a menos que se multiplique \(n\) por 4.

Se o tempo \(t\) passa de 1 para 3, por qual fator é preciso multiplicar \(n\) para manter o mesmo erro?

Se com \(n = 10\) passos o erro é \(0{,}04\), quantos passos são necessários para chegar a \(0{,}01\) (mesmo \(t\))?

O custo é claro: mais precisão pede circuitos mais fundos, e em hardware ruidoso (veja o curso de hardware) profundidade é cara. Fórmulas de ordem maior reduzem o erro por passo, mas exigem mais portas por passo.

\subsubsection{\texorpdfstring{O circuito de \(e^{-i\theta Z\otimes Z}\)}{O circuito de e\^{}\{-i\textbackslash theta Z\textbackslash otimes Z\}}}\label{o-circuito-de-e-itheta-zotimes-z}

Falta implementar cada termo. Para um termo de um qubit, é uma rotação (\(R_x\), \(R_y\) ou \(R_z\)). Para \(Z_1Z_2\), a exponencial \(e^{-i\theta Z\otimes Z}\) multiplica os estados da base por uma fase que depende da \textbf{paridade} (se os dois bits são iguais ou diferentes):

\begin{itemize}
\tightlist
\item
  \(\ket{00}\) e \(\ket{11}\) (paridade par, \(ZZ=+1\)): fase \(e^{-i\theta}\);
\item
  \(\ket{01}\) e \(\ket{10}\) (paridade ímpar, \(ZZ=-1\)): fase \(e^{+i\theta}\).
\end{itemize}

O truque é calcular a paridade no qubit 2 com uma CNOT, aplicar a fase com uma rotação \(R_z\) nele, e desfazer a CNOT:

\[
e^{-i\theta Z\otimes Z} = \text{CNOT}\,\big(I\otimes R_z(2\theta)\big)\,\text{CNOT}.
\]

Confira com \(\ket{01}\): a CNOT leva a \(\ket{01}\) (controle 0, nada muda), o qubit 2 está em \(\ket{1}\), onde \(R_z(2\theta)\) dá \(e^{+i\theta}\), e a segunda CNOT desfaz. Resultado: \(e^{i\theta}\ket{01}\), como previsto. Em \(\ket{11}\) a CNOT vira \(\ket{10}\), o alvo em \(\ket{0}\) ganha \(e^{-i\theta}\), e a segunda CNOT volta para \(\ket{11}\).

No circuito, a fase \(e^{-i\theta ZZ}\) de um estado com paridade ímpar é \(e^{+i\theta}\). Para \(\theta = 0{,}2\), qual é a fase (em radianos, com sinal) que o estado \(\ket{10}\) recebe em \(e^{-i\theta Z\otimes Z}\)?

Para termos como \(X_1X_2\) ou \(Y_1Y_2\), basta antes e depois girar a base (com \(H\) para \(X\); com \(S^\dagger\) e depois \(H\) para \(Y\)) e usar o mesmo circuito.

\subsection{Algoritmos variacionais e VQE}
Hoje os computadores quânticos têm dezenas a poucas centenas de qubits ruidosos e não corrigem erros. Dá para fazer algo útil com eles mesmo assim? A aposta dos algoritmos variacionais é: usar o chip quântico só para o que ele faz bem, deixar o resto para um computador comum, e tolerar circuitos curtos.

\subsubsection{A era NISQ}\label{a-era-nisq}

\textbf{NISQ} (do inglês \emph{noisy intermediate-scale quantum}) é o nome da fase atual: escala intermediária (algumas dezenas a centenas de qubits) e com ruído. Circuitos profundos, como o de estimação de fase com \(e^{-iHt}\) longos, acumulam erros demais. A saída é usar circuitos \textbf{rasos}, controlados por poucos números ajustáveis, e deixar um otimizador clássico procurar bons valores.

\subsubsection{Circuito parametrizado (ansatz)}\label{circuito-parametrizado-ansatz}

Um \textbf{ansatz} é uma família de circuitos com parâmetros reais \(\vec\theta = (\theta_1,\dots,\theta_m)\), que prepara o estado \(\ket{\psi(\vec\theta)}\). Exemplo de 1 qubit:

\[
\ket{\psi(\theta)} = R_y(\theta)\ket{0} = \cos\tfrac{\theta}{2}\ket{0} + \sin\tfrac{\theta}{2}\ket{1}.
\]

Variar \(\theta\) de \(0\) a \(2\pi\) percorre o círculo inteiro da esfera de Bloch no plano \(xz\): \(\ket{0}\) (\(\theta=0\)), \(\ket{+}\) (\(\pi/2\)), \(\ket{1}\) (\(\pi\)), \(\ket{-}\) (\(3\pi/2\)) e de volta. Neste texto os ângulos estão em radianos, salvo quando indicado. Para vários qubits, o ansatz típico alterna camadas de rotações com portas de emaranhamento (CNOT). Quanto mais \textbf{expressivo} o ansatz (isto é, quanto mais estados ele consegue produzir), mais perto do estado fundamental ele chega, mas mais fundo e mais difícil de otimizar ele fica.

\subsubsection{O princípio variacional}\label{o-princuxedpio-variacional}

Muitos problemas (química, materiais, otimização) pedem o \textbf{estado fundamental}: o estado de menor energia \(E_0\) de um hamiltoniano \(H\) (o menor autovalor). Para qualquer estado \(\ket{\psi}\) a \textbf{energia esperada} é

\[
E = \bra{\psi}H\ket{\psi}.
\]

O fato central: \textbf{nenhum estado tem energia esperada abaixo de \(E_0\)}. A razão: escreva \(\ket{\psi} = \sum_j c_j\ket{E_j}\) em autovetores de \(H\), com \(\sum|c_j|^2=1\). Então \(E = \sum_j |c_j|^2 E_j\) é uma média ponderada dos autovalores, e uma média nunca fica abaixo do menor valor. Logo, para o ansatz,

\[
E(\vec\theta) = \bra{\psi(\vec\theta)}H\ket{\psi(\vec\theta)} \;\ge\; E_0.
\]

Minimizar \(E(\vec\theta)\) dá uma \textbf{cota superior} de \(E_0\) que é exata se o ansatz contém o estado fundamental.

Exemplo: \(H = Z\) tem autovalores \(+1\) e \(-1\). O estado \(\ket{1}\) dá \(E=-1=E_0\). O estado \(\ket{+}\) dá \(E = 0\) (média de \(+1\) e \(-1\) com pesos iguais), acima de \(E_0\), como previsto.

Qual é o menor valor possível de \(\bra{\psi}H\ket{\psi}\) para \(H = Z + X\) (o menor autovalor)? Dica: \((Z+X)^2 = 2I\), então os autovalores são \(\pm\sqrt2\).

\subsubsection{\texorpdfstring{Medir \(\langle H\rangle\) por termos de Pauli}{Medir \textbackslash langle H\textbackslash rangle por termos de Pauli}}\label{medir-langle-hrangle-por-termos-de-pauli}

O chip só entrega bits de medida, não \(\bra\psi H\ket\psi\) diretamente. Mas \(H = \sum_k c_k P_k\) (veja a aula de trotterização), e pela linearidade

\[
\bra{\psi}H\ket{\psi} = \sum_k c_k \bra{\psi}P_k\ket{\psi}.
\]

Cada \(\langle P_k\rangle\) se estima repetindo o circuito muitas vezes (\emph{shots}) e medindo na base de \(P_k\): o resultado é \(+1\) ou \(-1\) e o valor esperado é a média. Para \(Z\), mede-se direto. Para \(X\), aplica-se antes uma porta \(H\) (Hadamard) para trocar de base e então mede-se em \(Z\). Para \(Y\), usa-se \(S^\dagger\) e depois \(H\).

Exemplo: se \(\langle Z\rangle = 0{,}6\) e \(\langle X\rangle = 0{,}8\), então para \(H = 0{,}5Z + 0{,}3X\) temos \(E = 0{,}5\cdot0{,}6 + 0{,}3\cdot0{,}8\).

Calcule esse valor.

Qual porta se aplica antes de medir em \(Z\) para estimar \(\langle X\rangle\)?

Como cada termo estimado vem de um número finito de shots, \(E(\vec\theta)\) tem ruído estatístico, que cai com \(1/\sqrt{\text{shots}}\). Reduzir o ruído por um fator 10 exige quantas vezes mais shots?

\subsubsection{O laço híbrido}\label{o-lauxe7o-huxedbrido}

O VQE (\emph{variational quantum eigensolver}) junta as peças em um laço:

\begin{enumerate}
\def\labelenumi{\arabic{enumi}.}
\tightlist
\item
  O computador \textbf{clássico} escolhe \(\vec\theta\).
\item
  O \textbf{chip quântico} prepara \(\ket{\psi(\vec\theta)}\) e mede os termos de \(H\).
\item
  O clássico soma tudo, obtém \(E(\vec\theta)\), e um \textbf{otimizador} (por exemplo descida de gradiente; COBYLA e Nelder-Mead são só nomes de otimizadores clássicos que não usam derivadas) propõe novos \(\vec\theta\).
\item
  Repete até a energia parar de cair.
\end{enumerate}

O quântico entrega só uma função \(E(\vec\theta)\); a otimização é clássica. Limitações reais: o ruído pode esconder o mínimo, a \textbf{paisagem} de \(E(\vec\theta)\) (o gráfico da energia em função dos parâmetros, com seus vales e planícies) pode ter muitos mínimos locais e regiões planas, e nada garante que o ansatz alcance o estado fundamental.

\subsubsection{Exemplo completo de 1 qubit}\label{exemplo-completo-de-1-qubit}

Tome \(H = Z + 0{,}5X\) e o ansatz \(R_y(\theta)\ket{0}\). Para esse estado, \(\langle Z\rangle = \cos\theta\) e \(\langle X\rangle = \sin\theta\), então

\[
E(\theta) = \cos\theta + 0{,}5\sin\theta.
\]

Em \(\theta=0\) (\(\ket{0}\)): \(E=1\). Em \(\theta = \pi/2\) (\(\ket{+}\)): \(E = 0{,}5\). O mínimo está em \(\theta = \pi + \arctan(0{,}5) \approx 3{,}60\) rad (\(206{,}6^\circ\)), fora do intervalo \([0,\pi]\), e vale \(-\sqrt{1+0{,}25} \approx -1{,}118\). Em geral, o mínimo de \(a\cos\theta + b\sin\theta\) é \(-\sqrt{a^2+b^2}\). Esse valor é exatamente o menor autovalor de \(H\): o ansatz de 1 qubit já cobre o estado fundamental.

Antes de otimizar: quanto vale \(E(\theta) = \cos\theta + \sin\theta\) em \(\theta = \pi/2\)?

Para \(H = Z + X\), em qual \(\theta\) (em graus, entre 0° e 360°) \(E(\theta) = \cos\theta + \sin\theta\) é mínima? (Dica: ache o ângulo em radianos e converta.)

\subsection{QAOA}
Três amigos querem formar dois times para um jogo, mas cada par de amigos que se desentende só atrapalha se ficar no mesmo time. Como separar o máximo de pares de desafetos? É um problema de otimização combinatória, e o QAOA ataca esse tipo de problema com um circuito parametrizado, no mesmo espírito do VQE.

\subsubsection{Otimização combinatória e MaxCut}\label{otimizauxe7uxe3o-combinatuxf3ria-e-maxcut}

Um problema \textbf{combinatório} pede a melhor entre muitas escolhas discretas (aqui, \(2^n\) divisões possíveis). \textbf{MaxCut}: dado um grafo (vértices ligados por arestas), divida os vértices em dois grupos para maximizar o número de arestas que ligam vértices de \textbf{grupos diferentes}. Essas arestas são as ``cortadas''.

Representamos a divisão por uma string de bits \(z = z_1 z_2\dots z_n\): o bit diz o grupo de cada vértice. O \textbf{valor do corte} \(C(z)\) é simplesmente o número de arestas cujos extremos têm bits diferentes.

Exemplo: um caminho \(0 - 1 - 2\) (arestas \((0,1)\) e \((1,2)\)) com \(z = 010\) separa o vértice 1 dos outros dois: as duas arestas são cortadas, \(C = 2\). Com \(z=000\), \(C=0\).

Qual é o valor máximo do corte de um triângulo (3 vértices, as 3 arestas)? Dica: com 2 grupos, dois vértices sempre ficam juntos.

\subsubsection{Do corte ao hamiltoniano}\label{do-corte-ao-hamiltoniano}

Para escrever \(C(z)\) como operador, use o sinal \(s_i = (-1)^{z_i}\) (\(+1\) para bit 0, \(-1\) para bit 1). Cada aresta contribui \(\frac{1 - s_is_j}{2}\): vale \(1\) se os sinais diferem e \(0\) se são iguais. Trocando \(s_i\) por \(Z_i\), o \textbf{hamiltoniano de custo} é

\[
H_C = \sum_{(i,j)} \frac{1 - Z_iZ_j}{2}.
\]

Ele é diagonal na base computacional: \(H_C\ket{z} = C(z)\ket{z}\). Checagem de uma aresta: em \(\ket{01}\) temos \(Z_0Z_1\ket{01} = (+1)(-1)\ket{01} = -\ket{01}\), logo \(\frac{1-(-1)}{2} = 1\): aresta cortada. Em \(\ket{00}\), \(Z_0Z_1=+1\) e o valor é \(0\).

Quanto vale \(\frac{1 - Z_0Z_1}{2}\) aplicado a \(\ket{01}\) (o autovalor)?

Maximizar o corte é achar o autoestado de \textbf{maior} autovalor de \(H_C\) (ou o de menor de \(-H_C\), tal como no VQE). Note que \(H_C\) é uma soma de termos \(ZZ\), exatamente os que a aula de trotterização implementa com CNOT-Rz-CNOT.

\subsubsection{O circuito do QAOA}\label{o-circuito-do-qaoa}

O \textbf{QAOA} (\emph{quantum approximate optimization algorithm}) é construído em três passos.

\textbf{Passo 1, estado inicial.} Uma Hadamard em cada qubit dá \(\ket{+}^{\otimes n}\), a superposição uniforme de todas as divisões, cada uma com probabilidade \(2^{-n}\).

\textbf{Passo 2, operador de custo.} Aplique \(e^{-i\gamma H_C}\), que multiplica cada \(\ket z\) pela fase \(e^{-i\gamma C(z)}\). Para cada aresta, é um \(e^{-i\theta ZZ}\) com \(\theta = -\gamma/2\) (o \(\tfrac12\) constante é fase global), feito com CNOT-Rz-CNOT. Divisões com cortes diferentes ganham fases diferentes, mas fases não mudam probabilidades.

Teste: depois de um operador de custo qualquer, partindo de \(\ket{++}\) (2 qubits), qual é a probabilidade de medir \(00\)?

\textbf{Passo 3, mixer.} Para converter as fases em probabilidades, aplique o \textbf{mixer} \(e^{-i\beta H_B}\), com \(H_B = \sum_i X_i\). Ele é um \(R_x(2\beta)\) em cada qubit e mistura as amplitudes: sem ele, o custo seria invisível na medida.

\textbf{Repetir.} Os passos 2 e 3 formam uma camada, com seus próprios \((\gamma_k,\beta_k)\). Com \(p\) camadas:

\[
\ket{\gamma,\beta} = \underbrace{e^{-i\beta_p H_B}\, e^{-i\gamma_p H_C}}_{\text{camada } p}\cdots \underbrace{e^{-i\beta_1 H_B}\, e^{-i\gamma_1 H_C}}_{\text{camada } 1}\ket{+}^{\otimes n}.
\]

São \(2p\) parâmetros no total.

Quantos parâmetros reais livres tem um QAOA com \(p = 3\) camadas?

Com \(\beta = \pi/8\), qual é o ângulo (em radianos) da porta \(R_x(2\beta)\) aplicada em cada qubit?

\subsubsection{O laço, de novo híbrido}\label{o-lauxe7o-de-novo-huxedbrido}

Como no VQE, o QAOA é um laço híbrido: o otimizador clássico escolhe \((\vec\gamma,\vec\beta)\); o chip prepara \(\ket{\gamma,\beta}\) e mede muitas vezes, estimando \(\langle H_C\rangle = \bra{\gamma,\beta}H_C\ket{\gamma,\beta}\), o corte médio. O otimizador procura parâmetros que maximizem esse valor. No fim, mede-se o estado e, entre as strings que saem, fica-se com a de maior corte (que é calculada de forma clássica).

Ponto de partida: no estado uniforme, cada aresta é cortada com probabilidade \(\tfrac12\) (os dois bits são independentes e diferentes em metade dos casos). Então \(\langle H_C\rangle = \tfrac12\cdot(\text{número de arestas})\).

Qual é o corte médio do estado uniforme no triângulo?

Como o máximo é 2, o QAOA já parte de \(1{,}5/2 = 75\%\) do ótimo; cada camada tenta aumentar esse número. Com os melhores parâmetros, aumentar \(p\) nunca piora o resultado e, no limite \(p\to\infty\), tende ao ótimo (na prática, achar esses parâmetros fica mais difícil).

\subsubsection{Ligação com evolução adiabática e trotterização (leitura opcional)}\label{ligauxe7uxe3o-com-evoluuxe7uxe3o-adiabuxe1tica-e-trotterizauxe7uxe3o-leitura-opcional}

Imagine um hamiltoniano que muda devagar de \(-H_B\) (estado fundamental \(\ket{+}^{\otimes n}\), o inicial) para \(-H_C\) (estado fundamental igual ao melhor corte). O \textbf{teorema adiabático} (um sistema que muda devagar o bastante fica no estado de menor energia) diz que, se a mudança é lenta o bastante, o sistema permanece no estado fundamental instantâneo e termina na resposta. Simular essa evolução com a fórmula de Trotter alterna passos de \(e^{-i\gamma H_C}\) e \(e^{-i\beta H_B}\), com \(\gamma\) e \(\beta\) pequenos: \textbf{é exatamente a estrutura do QAOA}. A diferença é que, no QAOA, os ângulos \(\gamma_k,\beta_k\) não são impostos pelo cronograma adiabático, e sim ajustados livremente pelo otimizador, o que permite obter bons resultados com poucas camadas.

